\begin{frame}{전체 그림: 어떻게 읽는가}
  \begin{itemize}
    \item 왼쪽$\to$오른쪽으로 \kb{데이터 $\to$ 파이프라인
      $\to$ 평가 $\to$ 발표} 로 흐름.
    \item \textbf{붉은 점선}이 핵심: test에서 \textbf{누수}
      (Week 6.3)나 \textbf{과적합}(Week 4.1)이 보이면
      돌아가서 다시.
    \item ``한 번 지나가면 돌아올 수 없는 것''은
      \textbf{test 한 번의 측정}뿐 --- 그 외 모든 단계는
      되돌아갈 수 있다.
  \end{itemize}
\end{frame}
